\section{Discussion session: open data (Day 1, late afternoon)}

\subsection{Open data proposal}
\speaker{Karolos Potamianos}{University of Warwick}
This contribution (no slides were posted on Indico) put forward a concrete
proposal, in the spirit of community datasets such as JetClass and the CMS Open
Data used by several speakers during the day, for a shared, well-documented
open dataset tailored to polarization tagging of hadronic boosted bosons:
common polarized signal samples ($W_L/W_T$, $Z_L/Z_T$) and QCD backgrounds
with agreed generator settings, truth-labeling conventions, and evaluation
metrics, so that the taggers presented at this workshop (ParT-based,
Lund-plane/QTTN, DDT-decorrelated, regression-based) can be compared on equal
footing and reused across ATLAS, CMS, and the phenomenology community.

\subsection{Open data roundtable}
\speaker{Antonio Giannini (USTC), Nurfikri Norjoharuddeen (LUT)}{conveners}
The roundtable gathered experimental and theory input on the practical aspects
of the proposal: which level of information to release (parton, particle,
Delphes-like, or full-simulation constituents), how to encode polarization
definitions consistently (frame choices, DPA/NWA conventions, treatment of
interference) so that samples remain theoretically meaningful, licensing and
hosting (Zenodo/CERN Open Data), and how a common benchmark could feed directly
into the COMETA deliverable discussed on Day 2. There was broad agreement that
a public, polarization-labeled boosted-boson dataset would significantly lower
the entry barrier for new groups and make performance claims comparable across
the field.
